\section*{Sets and Indices}

\[
P=\{\text{A1},\text{A2},\text{A3}\}
\]

Set of products, indexed by \(p \in P\).

\section*{Parameters}

For each product \(p \in P\):
\begin{itemize}
    \item \(D_p\): maximum demand (\texttt{Maximum\_Demand})
    \item \(s_p\): selling price per unit (\texttt{Selling\_Price})
    \item \(c_p\): production cost per unit (\texttt{Production\_Cost})
    \item \(q_p\): production quota in units per day (\texttt{Production\_Quota})
    \item \(F_p\): activation cost (\texttt{Activation\_Cost})
    \item \(B_p\): minimum batch size (\texttt{Minimum\_Batch})
\end{itemize}

Global parameter:
\begin{itemize}
    \item \(T\): available production days in the month (\texttt{production\_days})
\end{itemize}

Using the provided data for instance 22:
\[
T=22
\]
and
\[
\begin{array}{c|cccccc}
p & D_p & s_p & c_p & q_p & F_p & B_p \\
\hline
\text{A1} & 5300 & 124 & 73.30 & 500 & 170000 & 20 \\
\text{A2} & 4500 & 109 & 52.90 & 450 & 150000 & 20 \\
\text{A3} & 5400 & 115 & 65.40 & 550 & 100000 & 16
\end{array}
\]

\section*{Decision Variables}

For each product \(p \in P\):
\begin{itemize}
    \item \(x_p \ge 0\) (code: \texttt{x[p]}): production quantity of product \(p\)
    \item \(y_p \in \{0,1\}\) (code: \texttt{y[p]}): activation indicator for product \(p\)
\end{itemize}

\section*{Objective}

Maximize total profit:
\[
\max \sum_{p \in P} \left[(s_p-c_p)x_p - F_p y_p\right].
\]

\section*{Constraints}

The implemented model imposes that all three products are activated:
\[
y_p = 1 \qquad \forall p \in P.
\]

Demand upper bound linked to activation:
\[
x_p \le D_p\, y_p \qquad \forall p \in P.
\]

Minimum batch linked to activation:
\[
x_p \ge B_p\, y_p \qquad \forall p \in P.
\]

Total production time limit:
\[
\sum_{p \in P} \frac{x_p}{q_p} \le T.
\]

Variable domains:
\[
x_p \ge 0 \qquad \forall p \in P,
\]
\[
y_p \in \{0,1\} \qquad \forall p \in P.
\]

\section*{Notes on Implementation}

The formulation is implemented as a mixed-integer linear program in \texttt{current\_heuristic.py} and solved with Gurobi through PuLP.

Although binary activation variables \(y_p\) are declared, the code also explicitly enforces
\(
y_p=1
\)
for every product. Therefore, in the implemented model all products must be produced, activation costs are always incurred, and the lower bounds reduce to
\(
x_p \ge B_p
\)
while the upper bounds reduce to
\(
x_p \le D_p
\).

The production-days constraint uses the exact code expression
\[
\sum_{p\in P} x_p/q_p \le T,
\]
so each product consumes production time in proportion to its quantity divided by its quota. No integrality is imposed on \(x_p\); production quantities are continuous.
